\chapter{Compositional Propagation and Library Compatibility}
\label{ch:compositional}

A caller's specification is not independent of the libraries it calls. When a
method invokes another, the obligations the callee imposes become, after
substitution, obligations on the caller; the caller's contract is in part a
function of its callees' contracts. This dependence is the subject of the present
chapter, and it is the mechanism on which the thesis's adoption argument most
directly turns. If a caller's specification depends on its callees', then a change
to a library --- a new version that tightens a precondition or weakens a guarantee
--- propagates into the specifications of the code that uses it, and the
behavioural compatibility of an upgrade becomes a question answerable at the
specification level. The chapter measures how strong this dependence is on real
code, by implementing a top-down compositional analysis that propagates callee
specifications into callers, measuring what it recovers across five open-source
libraries totalling 21{,}052 methods, and characterising the structural shapes of
the recovered obligations --- which turn out to be exactly the shapes through which
a library change reaches its callers. The chapter answers RQ4, and connects the
result both backward to the test-generation value of Chapter~\ref{ch:testgen} and
forward to the compatibility application the dependence enables.

\section{Why the Dependence of Caller on Callee Matters}
\label{sec:comp-why}

The dependence of a caller's specification on its callees' is consequential for
two reasons, one bearing on the thesis's other results and one bearing on its
central adoption argument.

The first is that it sets a ceiling on those other results. The specifications the
earlier chapters validated, distributed, and supplied to a model were computed by
the inference engine's single-pass propagation, which captures only part of the
caller-on-callee dependence (Section~\ref{sec:comp-blind}). If a more thorough
analysis recovers substantially more of the dependence, then the specifications the
earlier chapters used were weaker than they could have been, and the value they
measured was a lower bound. How much the single pass leaves uncaptured therefore
determines whether the thesis's other results sit near the achievable ceiling or
well below it.

The second, and more important, reason is the compatibility application. The
behavioural compatibility of a library upgrade is, at bottom, a question about how
the library's changed contracts affect the contracts of the code that calls it ---
and that effect exists only to the extent that caller specifications depend on
callee specifications. If the dependence were weak, a library change would barely
register in its callers' contracts and specification-level compatibility analysis
would have little to work with; if it is strong, a library change propagates
visibly into caller contracts and the analysis has a rich signal. Measuring the
strength of the dependence is therefore the prerequisite for the compatibility
application: it establishes whether the signal that application would rely on is
substantial. The chapter finds that it is --- caller specifications depend on
callee specifications heavily, through cross-method obligations that a single-pass
analysis cannot even express --- which is what makes
Section~\ref{sec:comp-compatibility}'s account of compatibility analysis a
mechanism with real content rather than a hopeful extrapolation.

\section{Where Single-Pass Propagation Is Blind}
\label{sec:comp-blind}

The single-pass propagation of Section~\ref{sec:inf-inter} is the right
strategy for the simple case, and it is sound for that case: a precondition that
a callee requires \emph{unconditionally} must hold at the call site, and
conjoining the substituted precondition into the caller's set is correct. But
the strategy is structurally blind to two classes of cross-method obligation,
and the blindness is not a bug to be fixed within the single pass but a
consequence of the pass's architecture.

The first class is the \emph{branch-conditional} obligation. Consider a caller
whose body contains a call inside a conditional:

\begin{lstlisting}
if (g) {
    n(args);   // n requires p
}
\end{lstlisting}

The caller's precondition is not $p[\mathit{args}/\mathit{params}]$ but
$g \Rightarrow p[\mathit{args}/\mathit{params}]$ --- the obligation holds only
on the path where the call is reached. Single-pass propagation has two options,
both wrong. It can append the substituted precondition unconditionally, which
over-constrains the caller on every path where the call does not fire; or it can
skip the substitution to be safe, which misses the obligation entirely on the
path where the call does fire. Neither is correct, because producing the correct
clause requires the caller's control-flow context --- the guard $g$ enclosing
the call --- and a single substitution at the call site does not carry that
context.

The second class is the \emph{polymorphic-dispatch} obligation. Consider a call
through an interface or overridable method whose declared type has several
implementations. The runtime dispatch may select any of them, so a sound caller
precondition must hold for \emph{at least one} candidate: it is the disjunction
$p_1 \vee p_2 \vee \dots \vee p_k$ of the candidates' substituted preconditions.
Single-pass propagation resolves each call to a single nominal callee and cannot
produce this disjunction, because it never enumerates the candidate set.

It is worth being precise about why these two shapes, and not others, are the ones
the single pass misses. The single pass misses them because both require
information that lives in the \emph{caller}, not in the callee whose contract is
being propagated. The branch-conditional shape requires the guard enclosing the
call, which is a fact about the caller's control flow; the polymorphic shape
requires the set of types the receiver might hold, which is a fact about the
caller's context and the class hierarchy. Single-pass propagation, structured as a
substitution of the callee's contract at the call site, has the callee's contract
in hand but does not consult the caller's surrounding structure --- it reads the
callee and writes at the call site, and the caller's guards and the receiver's
possible types are exactly the surrounding structure it does not read. The shapes
the single pass \emph{does} capture --- unconditional preconditions on a
single-candidate call --- are precisely those that need nothing from the caller
beyond the call's arguments, which the substitution already has. The boundary
between what the single pass captures and what it misses is therefore not arbitrary:
it is the boundary between obligations determined by the callee and the call's
arguments, which a substitution handles, and obligations determined by the caller's
control-flow and type context, which a substitution cannot reach.

Both shapes are recoverable, but only by an analysis that does what the single
pass does not: walk the caller's body in a weakest-precondition fashion,
accumulating the guard context at each program point and enumerating the
dispatch candidates at each call site. The classical weakest-precondition
calculus~\cite{dijkstra_guarded_1975} addresses both structurally --- a call
site reached only on $g$ contributes $g \Rightarrow p$ by the conditional rule,
and a polymorphic call site contributes the disjunction over candidates by
construction. A complete WP transformer for arbitrary Java is out of reach, for
the reasons of Section~\ref{sec:bg-reasoning}, but a \emph{scaffold} that walks
the body in WP-flavoured order, lifts each call site's substituted callee
specification through its enclosing guards, and recognises polymorphic
candidates is implementable as a bounded pass over the AST. That scaffold is
what this chapter builds.

\section{The Refinement Pass}
\label{sec:comp-algorithm}

The pass takes a populated specification cache, a call graph over the same
methods, and a map from method signature to AST declaration, and it mutates the
cache in place by adding clauses no single pass could produce.

\subsection{Driver}

The driver iterates the strongly connected components of the call graph in
reverse topological order, so that callees are refined before callers. For a
singleton component --- a method not involved in recursion --- it refines the
method once. For a cyclic component --- a set of mutually recursive methods ---
it refines each method in the component and repeats until no method gains a new
clause or a fixpoint cap is reached, the cap bounding runtime on pathological
graphs and logging a warning when it fires. In practice the cap never fires on
the libraries measured: their call graphs are acyclic under the analyser's
signature-matching scheme, so each method is refined exactly once.

\subsection{Per-Method Refinement}

For each method, the refiner walks every call expression in the body and, at
each call site, performs five steps. First, it \emph{resolves the callee} by
matching the call's name against the cache. Second, it \emph{substitutes
arguments}, replacing each formal parameter in the callee's preconditions with
the actual argument expression; the substitution is regex-based with word
boundaries so that a formal name does not clobber a substring of an unrelated
identifier, and a side-effecting actual argument rejects the substitution as
unsound. Third, it \emph{lifts through enclosing guards}: it walks the AST
upward from the call, accumulating the conjunction of any conditional or ternary
guard whose body contains the call, and if guards are present the substituted
precondition becomes $\mathit{guard} \Rightarrow \mathit{substituted}$, with a
call in an else-branch contributing the negated guard. Fourth, it \emph{detects
polymorphic dispatch}: if the resolved callee has overriding implementations in
subtypes, found via the call graph's class-hierarchy index, it collects every
candidate's substituted preconditions and emits their disjunction, which is
sound because at least one candidate's contract must hold for the dispatch to be
safe. Fifth, it \emph{conjoins} the lifted clause into the caller's set if it is
not already present. The pass also propagates callee \emph{postconditions} into
the caller, gated on a termination flag: a postcondition of a possibly-
non-terminating callee describes a state established only if the callee returns,
so propagating it unconditionally would be unsound, and the pass propagates only
when the callee is known to terminate.

\subsection{Why Postcondition Propagation Needs a Termination Gate}

The postcondition propagation just described carries a soundness subtlety that the
precondition propagation does not, and it is worth isolating because it is the kind
of detail that distinguishes a sound compositional pass from a plausible-looking
but broken one. A precondition is an obligation the caller must discharge before
the call; propagating it is sound regardless of whether the callee terminates,
because the obligation is on the entry to the call. A postcondition is a guarantee
that holds \emph{after} the call returns; propagating it into the caller's contract
asserts that the caller's state, after the call, satisfies the callee's
postcondition --- which is true only if the call actually returned.

For a callee that always terminates this is unproblematic. For a callee that may
not terminate --- one with an unbounded loop or unbounded recursion --- the
postcondition describes a state reached only on the terminating executions, and
propagating it unconditionally would assert it on executions where the callee
diverges, which is unsound. The termination gate is conservative: when termination
is unknown the pass declines to propagate, missing a postcondition rather than
asserting an unsound one, consistent with the lower-bound character of the whole
measurement. The polymorphic case compounds the subtlety, the sound clause being a
disjunction over the terminating candidates' postconditions, handled identically to
the precondition disjunction. The gate is thus a discipline applied at every
postcondition propagation, and it is part of the reason the postcondition additions
of Table~\ref{tab:comp-headline} are smaller, on most libraries, than the
precondition additions.

\subsection{The Driver in Detail}

The driver's reverse-topological traversal of strongly connected components is
what makes the refinement well-founded, and the reasoning is worth setting out.
Refining a method means lifting its callees' contracts into it, which is only
sound if those callees' contracts are themselves complete; refining a caller
before its callee would lift an under-refined contract and miss clauses the
callee would later gain. Reverse-topological order over the condensation of the
call graph --- the directed acyclic graph of strongly connected components ---
guarantees that every callee outside a method's own component is fully refined
before the method is visited. Within a component, where methods are mutually
recursive and no acyclic order exists, the driver iterates: it refines every
method in the component, and if any gained a clause it refines them all again,
because a clause gained by one member may, on the next pass, lift into another.
The iteration is monotone --- clauses are only added, never removed --- and the
clause set per method is finite, so the iteration reaches a fixpoint; the
configurable cap exists only to bound runtime on pathological graphs and logs a
warning if it fires.

On the five libraries measured, the cap never fires, for an interesting reason:
under the analyser's signature-matching resolution, the call graphs are acyclic,
so every component is a singleton and every method is refined exactly once. The
fixpoint machinery is therefore dormant on the measured corpora, present for
soundness on graphs that would exercise it but never invoked. This is itself a
finding about real code: the mutual recursion that would force fixpoint iteration
is rare enough in the five libraries that the simpler singleton case dominates,
which is why the runtime cost (Section~\ref{sec:comp-empirical}) is so modest ---
the pass does one refinement per method and stops.

\subsection{Class-Hierarchy Analysis and Dispatch Precision}

The polymorphic-dispatch detection rests on a class-hierarchy index, and its
precision is governed by the same trade-off that has long characterised
call-graph construction for object-oriented languages. The simplest sound
over-approximation of a virtual call's targets is class-hierarchy analysis, which
takes the targets to be every override of the called method in any subtype of the
receiver's declared type. It is sound --- the runtime target is always among them
--- but imprecise, because it includes subtypes that the receiver could never
actually be at the call site. A more precise analysis, rapid-type analysis,
prunes the candidate set to subtypes that are actually instantiated somewhere in
the program, and flow-sensitive analyses prune further by tracking which types
reach the receiver.

The refinement pass uses the class-hierarchy index its call graph provides, which
is a class-hierarchy-analysis-level approximation: it enumerates the overrides of
the called method in the receiver type's subtypes and emits the disjunction over
their preconditions. The imprecision has a benign direction for the pass's
soundness: including a candidate the receiver could never be makes the disjunction
weaker --- easier to satisfy --- which cannot make the emitted precondition unsound,
only less informative. A disjunction over too many candidates is a precondition
that demands too little, which is safe; a disjunction that omitted a real candidate
would be a precondition that demands too much, which would be unsound, and the
hierarchy index avoids omission by including every override. The pass thus inherits
the standard class-hierarchy-analysis trade-off in its benign form: it may emit a
disjunction weaker than a more precise analysis would, but never one that is
unsound, and the disjunctive-shape counts of the measurement are therefore a sound
if conservative account of the polymorphic obligations the code imposes.

The prevalence of the disjunctive shape on Vavr, the library with the deepest
interface hierarchies, is in part a consequence of this approximation: a deep
hierarchy means many overrides, and class-hierarchy analysis enumerates them all,
so a call through a Vavr interface produces a disjunction over many candidates. A
more precise analysis would prune some, narrowing the disjunctions, but at the cost
of the flow-sensitive machinery the pass deliberately forgoes. The measurement
reports the class-hierarchy-analysis-level result, which is the sound conservative
one, and notes that a more precise dispatch analysis is a direction that would
sharpen the disjunctions without changing the qualitative finding that they are
present and substantial.

\subsection{What the Pass Is and Is Not}

The pass is a \emph{lifting} step over an existing single-pass result, not a
complete WP transformer. It does not compute the weakest precondition through
arithmetic or array statements --- it accepts the engine's clauses verbatim ---
it does not reason about aliasing, it does not lift through loops (the engine's
loop-invariant heuristics handle those in the single pass), and it does not
discharge clauses against a solver (the downstream OpenJML pipeline does that).
It does exactly enough --- branch-guard lifting and polymorphic-dispatch
disjunction --- to expose the two structural classes of additions that single-
pass propagation cannot produce. The measurement that follows asks whether those
additions are empirically substantial on real code, and the answer is the
chapter's contribution.

A question the structural measurement raises but does not by itself answer is
whether the refined specifications remain \emph{sound} --- whether the clauses the
pass adds discharge against the implementation the way the engine's other clauses
do. The pass is designed so that they should: the lifted clauses are the callees'
already-validated preconditions, substituted and guarded, and a sound callee
precondition lifted through a sound guard is itself sound. The substitution
disciplines of Section~\ref{sec:comp-substitution} preserve this --- they reject
exactly the substitutions that would break soundness --- so a clause the pass adds
is, in principle, dischargeable whenever the callee's original clause was. The
refined specifications therefore feed the same OpenJML validation layer the
engine's other output feeds, and a clause that fails to discharge is flagged like
any other, which means the second pass does not relax the soundness discipline of
the first. The measurement reports the clauses the pass adds, not their discharge
rate, because discharging the full refined corpus --- twenty thousand methods with
their refined contracts --- is a substantial verification effort deferred with the
downstream mutation evaluation as a completion of the study. The architectural
point stands regardless: the refined clauses are not a new class of unvalidated
output but lifted instances of already-validated callee clauses, passing through the
same validation gate.

\section{The Soundness of Substitution}
\label{sec:comp-substitution}

The refinement pass lifts a callee's precondition into the caller by substituting
actual arguments for formal parameters, and the soundness of that substitution is
not automatic --- it is a property the pass must actively preserve, and the cases
where it does not hold are the cases the pass must reject. The textbook WP rule
for procedure call substitutes actuals for formals, but the rule assumes the
actuals are pure expressions whose evaluation has no side effect and whose value
does not change between the precondition's evaluation point and the call. Real
Java arguments violate both assumptions, and the pass handles the violations by
refusing the unsound substitutions rather than performing them incorrectly.

The first hazard is a side-effecting actual. If a call passes \texttt{n(i++)} and
the callee requires its parameter be non-negative, naively substituting yields
\texttt{i++ >= 0}, a clause whose evaluation increments \texttt{i} and whose
meaning is therefore ill-defined as a precondition --- a precondition must be a
predicate the caller can evaluate without changing the state it constrains. The
pass detects increment, decrement, and assignment expressions in actual arguments
and rejects the substitution for those call sites, contributing no clause rather
than an unsound one. The conservatism costs recall --- a genuine obligation is
missed --- but it preserves the invariant that every clause the pass adds is one
the caller could in principle satisfy.

The second hazard is name capture in the substitution itself. Replacing a formal
parameter named \texttt{s} by a textual search-and-replace would corrupt an
unrelated identifier that contains \texttt{s} as a substring --- turning
\texttt{this.size} into \texttt{this.<arg>ize} --- which is a class of bug the
implementation history of the engine records as a real and recurring one. The
pass substitutes with word-boundary matching, so that only whole-identifier
occurrences of the formal are replaced, and a formal that happens to be a
substring of another identifier leaves that identifier untouched. The
word-boundary discipline is the textual analogue of capture-avoiding substitution
in the lambda calculus: it ensures that substituting for \texttt{s} affects
exactly the occurrences of the variable \texttt{s} and nothing else.

The third hazard concerns the guard lifting itself. A guard that is itself
side-effecting --- a condition that calls a mutating method --- cannot soundly be
lifted into a precondition, for the same reason a side-effecting actual cannot:
the precondition would change the state it describes. The pass rejects
side-effecting guards, lifting only guards that are pure boolean expressions over
the caller's state. Together these three disciplines --- reject side-effecting
actuals, substitute with word boundaries, reject side-effecting guards --- are
what make the lifted clauses sound additions to the caller's contract rather than
syntactically plausible but semantically broken strings, and they are the reason
the measurement of Section~\ref{sec:comp-empirical} reports a lower bound on the
compositional gain: every rejection is a clause the pass could have added under a
less careful policy but declines to add under a sound one.

\section{Empirical Measurement}
\label{sec:comp-empirical}

\subsection{Subjects and Method}

The pass is measured on the source jars of five open-source Java libraries
spanning four idiom families: Apache Commons Lang (utility/string code, 3{,}494
methods), Apache Commons IO (I/O abstractions, 2{,}111 methods), Apache Commons
Math (numerical algorithms, 6{,}645 methods), jOOL (functional-style
combinators, 3{,}691 methods), and Vavr (persistent collections, 5{,}111
methods). Combined, the corpus is 1{,}712 source files and 21{,}052 methods with
a non-empty isolated baseline specification. The breadth is deliberate: it spans
the imperative utility-class idiom that the earlier chapters emphasise and the
higher-order generic style of the functional libraries, so that the result
cannot be attributed to a single coding style. Commons Lang and IO are the same
libraries used in Chapters~\ref{ch:testgen} and~\ref{ch:distribution}, which
keeps cross-chapter comparisons free of corpus-difference confounds.

For each library, the harness parses every source file, builds the call graph,
runs the engine's standard inference (including its single-pass
\texttt{InterproceduralAnalyzer}) to produce the isolated baseline, snapshots the
per-method clause counts, runs the refinement pass, and diffs the post-refinement
specifications against the baseline, categorising each newly added precondition
by syntactic shape. All measurements are deterministic given the inputs, and each
is reported once.

\subsection{The Call Graph and the Measurement Harness}

The pass depends on a call graph, and the harness that takes the measurement
depends on resolving source methods to the cache, so both deserve description ---
not least because their limitations bound the measurement, as the threats section
records. The call graph is built by walking every method declaration and recording,
for each call expression in its body, the callee it resolves to, together with the
class-hierarchy relationships --- which classes extend which, which implement which
interfaces --- that the polymorphic-dispatch detection consults. The graph's nodes
are method signatures and its edges are calls; its condensation into strongly
connected components is what the driver iterates.

The resolution of a call to a callee is by signature matching, and here the harness
makes a deliberate simplification whose consequence the measurement must account
for. Two source methods of the same name and parameter count can, under erased
generics, share a bytecode descriptor, and the harness's name-and-arity matching
cannot always distinguish them. The consequence is that a residual of call sites
resolve to the wrong callee or fail to resolve, and the pass adds no clause for
those sites. Because a missed or misresolved call site can only cause the pass to
add \emph{fewer} clauses --- it never invents a clause from a wrong resolution, the
substitution being rejected when the arities disagree --- the effect is to
\emph{under}-count the compositional gain. Every number the measurement reports is
therefore a lower bound on the true gain a fully type-resolving harness would find,
which is the conservative direction for a claim that the gain is substantial: the
true gain is at least what is reported.

The measurement is otherwise deterministic. Given the source jars, the call graph,
the inference, and the refinement are all deterministic functions of the input, so
re-running the harness produces byte-identical output, which the threats section
records as verified across all five libraries. The harness snapshots the per-method
clause counts before refinement, runs the refinement, and diffs --- a procedure
that isolates exactly the clauses the second pass adds, attributing nothing to the
single pass that the single pass did not produce. This isolation is what licenses
the headline claim's form: not that the refined specifications are large, but that
they are larger than the single-pass baseline by a specific, measured amount, on
the same methods, under the same inference, with only the second pass varied.

\subsection{Headline Result (RQ4)}

Table~\ref{tab:comp-headline} reports the headline measurements across the five
libraries. The second pass strictly extends the isolated set on every library:
it adds at least one clause to between 24.33\% and 43.85\% of the methods in each
library's cache, a median of 36.33\%, and adds 42{,}370 new precondition clauses
and 25{,}203 new postcondition clauses in total. The single-pass propagation is
not subsuming on any library studied, and the gap is consistent across idiom
families.

\begin{table}[ht]
\centering
\small
\caption{Effect of compositional refinement over the isolated single-pass
baseline across five libraries. \emph{Refined} counts methods whose spec gained
at least one new clause; percentages are of methods in cache.}
\label{tab:comp-headline}
\begin{tabular}{lrrrrr}
\toprule
& \textbf{Lang} & \textbf{IO} & \textbf{Math} & \textbf{jOOL} & \textbf{Vavr} \\
\midrule
Methods in cache             & 3{,}494 & 2{,}111 & 6{,}645 & 3{,}691 & 5{,}111 \\
\midrule
Methods refined              & 1{,}532 & 767     & 2{,}492 & 898     & 1{,}799 \\
\hspace{1em}\emph{\% of cache} & 43.85\% & 36.33\% & 37.50\% & 24.33\% & 35.20\% \\
With new precondition & 1{,}160 & 522    & 1{,}929 & 802     & 1{,}225 \\
With new postcondition & 1{,}270 & 666   & 1{,}969 & 467     & 1{,}387 \\
\midrule
New precondition clauses     & 5{,}763  & 2{,}001 & 25{,}796 & 5{,}388 & 3{,}422 \\
\hspace{1em}\emph{avg/refined} & 4.97 & 3.83 & 13.37 & 6.72 & 2.79 \\
\hspace{1em}\emph{max on one} & 132   & 41 & 665 & 240 & 177 \\
New postcondition clauses    & 5{,}206  & 2{,}015 & 11{,}745 & 3{,}223 & 3{,}014 \\
\bottomrule
\end{tabular}
\end{table}

The fraction of methods refined varies with idiom in an interpretable way. The
utility library Commons Lang shows the highest fraction, because its delegating
helpers expose many guarded call sites; the functional library Vavr shows fewer
refined methods but concentrates its additions in disjunctive shapes from
polymorphic dispatch, a consequence of its deep interface inheritance. The
per-method dispersion is heavy-tailed --- a small number of large fa\c{c}ade
methods account for the maxima of 132 to 665 new clauses --- but the mean new
preconditions per refined method (2.79 to 13.37) sits well below those maxima on
every library, establishing that the refinement is a population-level effect
rather than an artefact of a handful of large methods.

\subsection{Shape Distribution}

Table~\ref{tab:comp-shapes} categorises the newly added preconditions by
syntactic shape. Branch-conditional implications --- the shape the single pass
structurally cannot produce, because it lacks the control-flow context the WP
walk supplies --- are the largest single shape on four of the five libraries,
peaking at 64.1\% on Commons Math, where heavy algorithmic code paths gate calls
on numerical preconditions. The exception is jOOL, where null checks dominate
because the higher-order combinators delegate to user-supplied lambdas through a
small set of pinch-point validation sites; even there, branch-conditional shapes
account for a substantial third of the additions. Disjunctive clauses --- the
polymorphic-dispatch signal, the other shape the single pass cannot produce ---
are present on every library and peak at 9.7\% on Vavr, the library with the
deepest interface inheritance chains.

\begin{table}[ht]
\centering
\small
\caption{Syntactic shape of the newly added preconditions across the five
libraries. Branch-conditional implications dominate on four of five; disjunctive
clauses are present on every library.}
\label{tab:comp-shapes}
\begin{tabular}{lrrrrr}
\toprule
\textbf{Shape (\% of new precond.)} & \textbf{Lang} & \textbf{IO} & \textbf{Math} & \textbf{jOOL} & \textbf{Vavr} \\
\midrule
Branch-conditional (\texttt{==>})  & 52.6\% & 42.8\% & 64.1\% & 34.1\% & 54.2\% \\
Null check                         & 21.8\% & 34.7\% & 24.1\% & 60.2\% & 30.9\% \\
Range / bound                      & 19.8\% & 18.9\% & 7.0\%  & 0.5\%  & 4.5\%  \\
Disjunctive (\texttt{||})          & 4.3\%  & 1.8\%  & 3.7\%  & 5.1\%  & 9.7\%  \\
Other                              & 1.5\%  & 1.7\%  & 1.2\%  & 0.2\%  & 0.6\%  \\
\bottomrule
\end{tabular}
\end{table}

Null and range/bound clauses make up most of the remainder. These are
obligations the single pass \emph{would} produce when the call site is
unguarded; the second pass adds them when the call site is guarded by a
condition the single pass would not lift. The dominance of branch-conditional
and disjunctive shapes is the chapter's central evidence: between 38\% and 74\%
of the new clauses on each library are of shapes that the single pass cannot
express by construction, which is direct confirmation that the architectural
distinction between single-pass propagation and compositional refinement is real
and consequential, not merely a difference in implementation thoroughness.

\subsection{Interpreting the Magnitude}

The headline figure --- 42{,}370 new precondition clauses --- is large enough to
warrant interpretation, because a raw count can mislead in either direction. On one
hand, the count aggregates over twenty-one thousand methods, so the per-method
average is modest, and a reader should not imagine every method gaining hundreds of
clauses; the median refined method gains a handful, and the large counts are
concentrated in the fa\c{c}ade methods that fan out to many callees. On the other
hand, the count is a lower bound, for the reasons the substitution-soundness and
callee-resolution discussions establish: every rejected substitution and every
unresolved call site is a clause the pass declined to add under a sound policy, so
the true compositional gain exceeds the reported one.

The figure that matters more than the raw count is the fraction of methods refined
--- between a quarter and nearly a half of each library's cache --- because it
speaks to how pervasive the phenomenon is. That single-pass propagation leaves
something to recover on a quarter to a half of all methods, across five libraries
of different idioms, establishes that the gap is not a corner case affecting a few
unusual methods but a routine consequence of how cross-method obligations arise. A
gap affecting one method in a hundred would be a curiosity; a gap affecting one in
three is a structural property of the single-pass architecture, and it is the
pervasiveness, not the raw clause count, that makes the case for the second pass.

The shape distribution sharpens the interpretation once more. If the additions were
dominated by null checks and range bounds --- the shapes the single pass produces
when call sites are unguarded --- one could argue the second pass merely recovers
clauses the single pass missed by oversight, and a more careful single pass would
close the gap. But the additions are dominated by branch-conditional implications
and include disjunctions on every library, and these are shapes the single pass
cannot produce by construction, not by oversight. The magnitude is therefore
significant not merely for its size but for its composition: it is large, it is
pervasive, and it is concentrated in exactly the shapes that demonstrate the
architectural distinction is real.

\subsection{Runtime Cost}

The pass completes in 44 seconds total across all five libraries and 21{,}052
methods, an average of 2.1 milliseconds per method, ranging from 0.7 seconds on
Commons IO to 20.3 seconds on the largest library, Commons Math. No strongly
connected component hit the fixpoint cap on any library, because the call graphs
are acyclic under the analyser's signature-matching scheme and each method is
refined once. The cost is dominated by the per-method AST walk for call sites,
which is linear in body size; the lifting and substitution steps add constant
factors per call site. The refinement is therefore an addition of seconds to
tens of seconds on top of an inference pass that already takes minutes on
libraries of this scale --- a clearly affordable enhancement, which answers the
cost half of RQ4.

The scaling behaviour across the five libraries confirms the linear character. The
twenty-second figure on Commons Math, the largest library, against the sub-second
figure on Commons IO, the smallest, tracks the ratio of their method counts and
call-site densities rather than any superlinear effect: Math has both more methods
and, as a numerical library, more call sites per method, and its runtime reflects
the product. No library exhibits the runtime explosion that a fixpoint iteration on
a deeply cyclic call graph would produce, because no library's call graph is
cyclic under the analyser's resolution, and the dormant fixpoint machinery
contributes nothing to the measured times. The practical consequence is that a
deployer can predict the refinement cost from the corpus size --- roughly two
milliseconds per method --- without worrying about pathological cases, because the
pathological case (deep mutual recursion forcing many fixpoint iterations) does not
arise in the libraries measured and would, if it did, be bounded by the
configurable cap. The cost is thus not only affordable but predictable, which is
the property a deployer integrating the pass into a build needs.

\section{The Character of Each Library}
\label{sec:comp-perlibrary}

The five libraries were chosen to span idiom families, and the shape
distributions of Table~\ref{tab:comp-shapes} reflect their characters in ways
that illuminate what drives compositional refinement. Reading the libraries one
at a time turns the aggregate result into an account of \emph{why} the second
pass finds what it finds.

Apache Commons Lang, the utility library, shows the highest fraction of methods
refined (43.85\%) and a branch-conditional-dominated shape distribution. This
follows from its design: it is a collection of small static helpers that delegate
to one another and to validation utilities, and the delegation happens inside
guards --- a helper checks a condition and, on one branch, calls a more primitive
helper that imposes its own requirement. The guarded delegation is exactly the
structure the branch-conditional lifting recovers, and the library's density of
it is why nearly half its methods gain a clause. The library is, in a sense, the
ideal case for the second pass: small methods, heavy delegation, guarded calls.

Commons IO, the I/O abstraction library, shows a lower refined fraction (36.33\%)
and a higher proportion of null checks among its additions. Its methods are
larger and its delegation shallower, but its file and stream abstractions impose
many non-null requirements --- a stream must be open, a path must be non-null ---
that propagate through call chains. The branch-conditional fraction is still
substantial, reflecting the defensive guarding common in I/O code, but the null
check is more prominent than in Lang because the domain's central obligations are
non-nullness obligations.

Commons Math, the numerical library, is the outlier in magnitude: it adds 25{,}796
new precondition clauses, more than the other four libraries combined, and its
branch-conditional fraction peaks at 64.1\%. The explanation is the structure of
numerical code, which is dense with guarded calls into validation and computation
helpers --- a method computing a statistic checks the data's validity and, on the
valid branch, calls a helper that requires it, and the helper in turn guards
further calls. The deep, guarded call chains characteristic of numerical
algorithms produce branch-conditional obligations at every level, and the second
pass lifts them all, which is why a single large fa\c{c}ade method can gain
hundreds of clauses. Commons Math is the case that most stresses the pass and most
rewards it.

jOOL, the functional-combinator library, inverts the pattern: null checks
dominate its additions at 60.2\%, and branch-conditional clauses, while still a
substantial third, are less prevalent than elsewhere. The reason is the
higher-order style. jOOL's combinators take user-supplied functions and delegate
to them through a small set of validation pinch-points --- a combinator checks
its function arguments are non-null and then applies them --- so the dominant
propagated obligation is the non-nullness of the supplied function, not a guarded
numerical precondition. The functional idiom concentrates obligations at the
boundary where user code enters, and that boundary is a null check.

Vavr, the persistent-collection library, shows the highest disjunctive fraction
at 9.7\%, the polymorphic-dispatch signal. Its deep interface inheritance --- a
value type implementing a traversable interface implementing an iterable
interface, and so on --- means that a call through one of its interfaces may
dispatch to any of many implementations, and the second pass emits the
disjunction over candidates at each such site. Vavr is the case that most
exercises the polymorphic-dispatch half of the pass, and its disjunctive fraction
is direct evidence that the half is doing real work: a library that uses
polymorphism heavily produces the disjunctive obligations a library that does not
would lack.

The five libraries thus form a spectrum, and the second pass responds to each
according to its idiom: guarded delegation in Lang and Math produces
branch-conditional clauses, boundary validation in IO and jOOL produces null
checks, and deep polymorphism in Vavr produces disjunctions. The aggregate result
--- that the single pass is not subsuming on any of them --- holds across the
spectrum, but the \emph{reason} it holds differs by library, and the per-library
reading is what shows that the result is not an artefact of one coding style but a
general consequence of how cross-method obligations arise in code.

\section{From Compositional Dependence to Compatibility Analysis}
\label{sec:comp-compatibility}

The measurement establishes the fact the compatibility application rests on: a
caller's specification depends substantially on its callees', through cross-method
obligations a single-pass analysis cannot express. This section develops what that
fact makes possible, which is the thesis's most direct answer to the adoption
question --- a capability manual specification never made routine.

Consider a consumer whose code calls a library, and a new release of that library.
Three kinds of change can occur at the library boundary. The library may change a
method's \emph{signature}, which the compiler detects; this is the only kind of
incompatibility today's tooling reliably catches. It may preserve the signature but
change the method's \emph{behaviour} --- tightening a precondition so that inputs
the old version accepted now violate the contract, or weakening a postcondition so
that guarantees the caller relied on no longer hold. This second kind, behavioural
incompatibility under signature compatibility, is the dangerous one: it compiles,
it may even pass shallow tests, and it breaks in the cases the tests do not cover.
The third kind is no change at all, and a consumer wants to distinguish a
genuinely compatible upgrade from a dangerous one without running their entire
test suite against every release and hoping the suite covers the changed
behaviour.

Inferred, library-boundary specifications make the second kind detectable. The
inferrer, run over the old and the new library versions, produces a specification
for each method in each version; comparing the two reveals which methods' contracts
changed and how --- which preconditions strengthened, which postconditions weakened.
A strengthened callee precondition is a candidate incompatibility: the caller may
pass arguments the new version rejects. A weakened callee postcondition is another:
the caller may rely on a guarantee the new version no longer makes. And because the
compositional analysis of this chapter propagates callee contracts into caller
contracts, the analysis can go further than flagging the library-side change: it
can recompute the \emph{caller's} specification against the new library and compare
it to the caller's specification against the old, surfacing precisely the caller
obligations that the library change altered. A caller obligation that appears only
against the new version is a new demand the upgrade places on the caller; one that
disappears is a guarantee the caller has lost. The cross-method obligations this
chapter measured --- the guard-conditional clauses lifted from guarded calls, the
disjunctions from polymorphic dispatch --- are exactly the channels through which a
library change reaches the caller, which is why the shape distribution matters: it
shows that the propagation carries the kinds of obligation that behavioural change
travels through.

The thesis establishes the mechanism --- the substantial, measured dependence of
caller specifications on callee specifications --- but does not run the
version-to-version experiment that would demonstrate compatibility detection
end to end. That experiment, recomputing caller specifications against successive
library versions and reporting the incompatibilities the comparison surfaces, is
the foremost item of future work (Chapter~\ref{ch:conclusion}), and the present
chapter's contribution is to show that the signal it would rely on is real and
substantial. What the chapter demonstrates is the prerequisite: that a library's
specifications propagate into its callers' strongly enough, and through the right
structural channels, for their version-to-version change to carry the behavioural-
compatibility information a consumer needs. The compatibility application is the
capability this dependence enables, and naming it precisely --- mechanism
demonstrated here, end-to-end experiment deferred --- is the honest statement of
how far the thesis takes it.

\section{Preliminary Downstream Signal}
\label{sec:comp-downstream}

The structural result establishes that the second pass adds clauses; it does not
by itself establish that the added clauses help a downstream consumer. A
preliminary probe connects the refinement back to the test-generation result of
Chapter~\ref{ch:testgen}. The experiment re-runs the test-generation comparison
with two configurations of the specification context: the isolated baseline
specification (the P3 condition of Chapter~\ref{ch:testgen}) and the
compositionally refined specification, on the eleven-class Commons Lang fa\c{c}ade,
five generation runs per class.

On the eight of eleven classes that returned extractable tests in both
configurations, the refined specification yields a 9.1\% higher mean test count,
with six of the eight showing an increase. The largest gains are on
validation-heavy classes --- \texttt{Validate} (+54.8\%) and \texttt{BooleanUtils}
(+27.5\%) --- where the refined specification's branch-conditional clauses give
the model additional case structure to test. The single substantial regression
is \texttt{CharUtils} ($-23\%$), where the disjunctive obligations from the
class's many overloaded predicate methods appear to be interpreted by the model
as one combined assertion rather than several distinct cases. The aggregate is a
weak positive signal with high per-class variance, and the small sample
precludes a meaningful significance test.

The \texttt{CharUtils} regression repays closer attention, because a negative
result in a downstream probe is as informative as a positive one. \texttt{CharUtils}
is dense with overloaded predicate methods --- variants of \texttt{isAscii} and its
relatives --- and the refinement pass, encountering calls through these overloads,
emits disjunctions over the candidates' preconditions. The disjunction is sound and
informative to a verifier, but to a language model generating tests it appears to
read as a single combined assertion rather than as an enumeration of cases, so the
model produces one test where the baseline specification's simpler clauses prompted
several. The richer specification, in this instance, gives the model a clause whose
structure it handles less well than the simpler clause it replaces. This is a
genuine finding about the interaction between specification shape and model
consumption: not every structurally richer specification is a better prompt, and a
disjunction that helps a verifier may hinder a generator that does not decompose it.
The observation does not undermine the aggregate positive signal --- six of eight
classes improve --- but it identifies a specific mechanism by which refinement can
backfire for a particular consumer, which is exactly the kind of nuance a downstream
evaluation should surface and a structural measurement alone would miss.

The signal is reported with its limitations stated plainly. It measures the
number of tests the model emits, not the mutation-killing rate, because the
generated tests did not all compile cleanly in the measurement environment ---
a tooling issue in the test-generation harness and in the model's occasional
emission of non-Java tokens, not a property of the specifications. A full
mutation-coverage comparison between baseline and refined specifications is the
natural next step and is identified as the top-priority follow-up. What the
probe establishes is narrower but real: the delta between baseline and refined
specifications is non-zero at the test-generation step, which is consistent with
the hypothesis that the additional clauses carry downstream value and motivates
the fuller evaluation.

The fuller evaluation, when it is done, should measure mutation score rather than
test count, on a corpus large enough to support a significance test, with the
generated tests compiling and running so that their fault-detection power can be
assessed directly. The hypothesis it would test is specific: that the
branch-conditional and disjunctive clauses the refinement adds equip a model to
generate tests that kill mutants the baseline specification's tests miss --- mutants
on the guarded paths the branch-conditional clauses describe and the dispatch
targets the disjunctions enumerate. The probe's positive-but-noisy test-count
signal is consistent with this hypothesis but does not establish it, because a test
count says nothing about whether the additional tests kill additional mutants. The
\texttt{CharUtils} regression is a caution that the relationship may not be uniform
--- a richer specification can, for a particular consumer and a particular clause
shape, produce fewer or weaker tests --- so the fuller evaluation should report the
per-class distribution, not only the aggregate. The probe's role in the thesis is to
establish that the question is worth the fuller study's cost: the delta is non-zero,
the mechanism is plausible, and the regression shows the answer is not trivially
positive, which is exactly the situation that warrants a careful measurement rather
than an assumption.

\section{Worked Examples of the Two Shapes}
\label{sec:comp-examples}

The two recoverable shapes are best seen on concrete code. Consider first a
caller that reaches a guarded call:

\begin{lstlisting}
// callee: requires divisor != 0
static int safeDivide(int a, int divisor) { return a / divisor; }

static int scale(int value, int factor) {
    if (factor != 0) {
        return safeDivide(value, factor);
    }
    return value;
}
\end{lstlisting}

The callee \texttt{safeDivide} carries the inferred precondition
\texttt{divisor != 0}. Single-pass propagation, encountering the call
\texttt{safeDivide(value, factor)}, substitutes the actual argument to obtain
\texttt{factor != 0} and faces the dilemma of Section~\ref{sec:comp-blind}:
adding it unconditionally would forbid calling \texttt{scale} with
\texttt{factor == 0}, which the method explicitly handles by returning
\texttt{value}; skipping it would lose the obligation. The refinement pass walks
the body, finds the call enclosed in the guard \texttt{factor != 0}, and lifts
the substituted precondition through the guard to obtain
\texttt{factor != 0 ==> factor != 0} --- which simplifies to a tautology here,
correctly reflecting that the guard already discharges the callee's requirement.
The example is deliberately degenerate to show the mechanism; in the common case
the guard and the callee's requirement differ, as when a call to a method
requiring a non-empty list is guarded by a size check, and the lifted clause
\texttt{size > 0 ==> list.size() > 0} carries genuine, non-tautological
information that the single pass cannot express.

Now consider polymorphic dispatch:

\begin{lstlisting}
interface Shape { double area(); }   // no precondition
class Disk implements Shape {
    // requires radius >= 0
    public double area() { return Math.PI * radius * radius; }
}
class Square implements Shape {
    // requires side >= 0
    public double area() { return side * side; }
}

static double totalArea(Shape s) { return s.area(); }
\end{lstlisting}

The call \texttt{s.area()} in \texttt{totalArea} dispatches at runtime to
whichever implementation \texttt{s} actually is. Single-pass propagation resolves
the call to the nominal type \texttt{Shape}, whose \texttt{area} declares no
precondition, and so adds nothing. The refinement pass, consulting the call
graph's class-hierarchy index, enumerates the implementing candidates
\texttt{Disk} and \texttt{Square}, collects their substituted preconditions, and
emits the disjunction \texttt{radius >= 0 || side >= 0}. The disjunction is the
sound caller obligation: for the dispatch to be safe, at least one candidate's
precondition must hold, namely the one corresponding to the runtime type. This is
the polymorphic-dispatch shape, and it is structurally invisible to a single pass
that resolves each call to a single nominal callee. The shape's prevalence on
Vavr (Table~\ref{tab:comp-shapes}), the library with the deepest interface
inheritance, is a direct consequence of how much polymorphic dispatch its design
employs.

\section{Positioning}
\label{sec:comp-positioning}

The unified review of compositional and interprocedural analysis
(Section~\ref{sec:bg-compositional}) located the position this propagation
occupies along three axes --- whether a technique produces JML, reasons
compositionally across procedure boundaries, and checks its reasoning against a
solver --- and found the cell it fills (produces JML, reasons compositionally by
weakest-precondition lifting, feeds an existing JML verifier) unoccupied by prior
work. The measurement of this chapter substantiates that positioning: it
demonstrates the compositional reasoning concretely, on real code at scale,
through the guard-lifting and dispatch-disjunction steps that distinguish the
approach from the genuinely-compositional memory-safety analysers (which do not
produce JML), the deductive verifiers (which consume rather than produce
specifications), and the template-and-prune inferrers (which do not lift callee
contracts into callers). The novelty is the combination, and the value the
chapter adds to the positioning is the evidence that the combination produces a
substantial, structurally characterisable body of obligations --- the evidence on
which the compatibility application of Section~\ref{sec:comp-compatibility} rests.

\section{Implications for Inferrer Architecture}
\label{sec:comp-implications}

The measurement settles a design question that the single-pass-versus-second-pass
choice poses for any specification inferrer, and the settlement has consequences
beyond the particular engine of this thesis. The question is whether the
additional architectural complexity of a compositional pass is justified by the
information it recovers, and the answer the measurement gives is conditional on
what the inferrer is for.

For an inferrer whose output feeds a consumer that benefits from richer
specifications --- a verifier, a test generator, a language model --- the second
pass is clearly worth integrating: it adds tens of thousands of clauses across the
corpus at a cost of seconds, and the clauses are of shapes the single pass cannot
produce. The cost-benefit is lopsided in the second pass's favour, and the
preliminary downstream signal suggests the added clauses carry value to exactly
the kind of consumer that motivates the thesis. For an inferrer whose output is
consumed eagerly and superficially --- an IDE showing hover text, a tool that
displays the first precondition --- the additional clauses may not justify their
runtime, and the second pass is reasonably left off. The thesis's design makes the
pass opt-in for precisely this reason: the right choice depends on the consumer,
and the measurement gives a deployer the information to choose.

A subtler implication concerns inferrers that wish to remain single-pass for
engineering simplicity. The measurement shows that the branch-conditional shape,
which dominates the gap on four of five libraries, requires only the caller's
local control-flow context --- the guard enclosing each call --- which a
single-pass analyser already traverses during its AST walk. Such an analyser could
lift its substituted clauses through enclosing guards before adding them,
capturing the dominant shape of the gap without the full machinery of a second
pass over the call graph. The polymorphic-dispatch shape is harder to capture in a
single pass, because it requires consulting the class hierarchy at the point of
propagation, which is the kind of cross-cutting lookup a second pass is naturally
organised to perform. The measurement thus decomposes the gap into an easy half a
single pass could capture with a local extension and a hard half that wants a
second pass, and it quantifies each, which is the information an inferrer designer
needs to decide how far to go.

\section{Threats to Validity}
\label{sec:comp-threats}

\paragraph{Internal validity.} The clause-count diffs are deterministic;
rerunning the harness produces byte-identical output, verified across all five
libraries. The shape categorisation uses regex-based bucketing with a documented
order of precedence, so the buckets are reproducible; a small residual of
arithmetic comparisons the regex does not match falls into the \emph{other}
bucket and does not affect the headline result.

\paragraph{External validity.} The five libraries deliberately span four idiom
families, including generics-heavy and numerical code, to address the concern
that the result is utility-class-specific; it holds across all five. Application
and framework code (web frameworks, middleware) have different call-graph shapes
and are not represented; generalisation to them is open. Google Guava, included
in the earlier chapters, is excluded here after a stack-overflow fix in the
call-graph builder's interface-inheritance walk, because its source-jar size
exceeded the measurement budget; the qualitative conclusion would not change but
the run was not completed.

\paragraph{Construct validity.} ``Refined'' is operationalised as gaining at
least one clause not in the baseline, using exact string equality, so two
semantically equivalent but syntactically different clauses count as distinct ---
which makes the new-clause count an upper bound on genuinely new logical
obligations, though the shape categorisation is unaffected. The callee-resolution
step matches on name and arity, so a residual of call sites is lost to overload
collisions; a missed call site can only under-count additions, so the reported
numbers are a lower bound on the true compositional gain on every library. The
downstream signal measures test count, not mutation kills, and rests on a sample
of eight classes; it is reported as preliminary for exactly these reasons.

\paragraph{Conclusion validity.} The structural claims are exact measurements,
not statistical inferences, so no significance test applies to them; the
conclusions hold for the corpora and the specific refinement implementation
reported. The qualitative finding --- that single-pass propagation cannot lift
through guards or dispatch, and that a second pass recovers substantial clauses
of exactly those shapes --- has a structural explanation that suggests it will
transfer, but quantitative transfer to other libraries is an open question.

\section{Summary}
\label{sec:comp-summary}

This chapter asked whether the single-pass interprocedural propagation that makes
inference practical captures the cross-method specification information a genuinely
compositional analysis would, and answered by building and measuring a top-down
weakest-precondition refinement pass. The answer is that the single pass is not
subsuming: across five libraries spanning four idiom families and 21{,}052 methods,
the second pass adds 42{,}370 precondition clauses, refining a quarter to nearly a
half of each library's methods, at a cost of seconds to tens of seconds. The
additions are dominated by branch-conditional implications and include
polymorphic-dispatch disjunctions on every library --- the two shapes the single
pass cannot produce, because both require the caller's control-flow or
class-hierarchy context that a substitution at the call site does not carry. The
per-library analysis showed the result holds across idioms for idiom-specific
reasons, and a preliminary downstream probe connected the additions back to the
test-generation value of Chapter~\ref{ch:testgen}, finding a positive but noisy
signal that motivates a fuller mutation-coverage evaluation.

The chapter's contribution is twofold: a measurement, at a scale and across a
breadth the literature has not previously reported, of what single-pass inference
leaves uncaptured; and a characterisation of that residual as two structurally
identifiable shapes rather than an undifferentiated gap. The measurement gives an
inferrer designer the evidence to decide whether a second pass is worth its
complexity, and the characterisation shows that the easy half of the gap --- the
branch-conditional shape --- could be captured by a local extension to a single
pass, while the harder half --- the disjunctive shape --- wants the second pass's
cross-cutting view of the class hierarchy. Beyond the design guidance, the
measurement supplies the prerequisite for the thesis's compatibility argument: it
establishes that a caller's specification depends on its callees' substantially and
through the structural channels --- guard-conditional and dispatch-dependent
obligations --- through which a library change propagates to the code that uses it.
This dependence is what makes specifications at the library boundary a basis for
behavioural-compatibility analysis, the capability that
Section~\ref{sec:comp-compatibility} develops and that the thesis offers as a
direct answer to the adoption question: a thing practitioners could not previously
do, made possible by inferred, propagating, library-level specifications. The
end-to-end version-comparison experiment that would demonstrate it remains future
work, but its mechanism is measured and real.
